\section{Modelling Methodology}
\label{sec:method}

\subsection{Model family budget and pipelines}

The brief allows two to three families and asks for depth rather than breadth. We
use three, reused across both task types, each arranged as a \emph{complexity
ladder} whose rung 0 is the family's simplest variant --- which is the only way
``the added complexity was worth it'' becomes a finding rather than an
assumption. \textbf{A --- Linear} runs \code{LinearRegression} /
\code{LogisticRegression} at rung 0 and \code{Ridge} / $\ell_2$-penalised
logistic at rung 1. \textbf{B --- Tree-based} runs \code{DecisionTree} $\to$
\code{RandomForest} $\to$ \code{HistGradientBoosting}. \textbf{C --- Ensemble} is
soft \code{Voting} and \code{Stacking} assembled from the winners of A and B; its
members are complete pipelines rather than bare estimators, because the two
families need different encodings and one shared preprocessor would force a
single choice on both. The full specification, search spaces and selected values
are in Appendix~\ref{app:hyperparameters}.

\paragraph{Why these rungs, from the literature rather than from our results.}
The Milestone 1 feedback observed that Phase 1 cited the original algorithm
papers and compared families only broadly, so the rungs are argued here before
any result is shown. Within family B they are the three ways the literature
knows to improve on one tree: a random forest bags decorrelated trees to cut
variance \citep{breiman2001rf}; boosting fits trees to residuals with shrinkage
to cut bias \citep{friedman2001gbm}. Across the large empirical comparisons --
ten algorithms on eleven problems \citep{caruana2006empirical}, 179 classifiers
on 121 datasets \citep{fernandezdelgado2014hundreds}, 45 tabular benchmarks
against deep architectures \citep{grinsztajn2022tabular, shwartzziv2022tabular}
and 41 classifiers on business credit data \citep{lessmann2015benchmarking} --
boosted and bagged tree ensembles are the consistent top tier, boosting usually
ahead by a small margin. The boosted learner is \code{HistGradientBoosting}
rather than a third-party library because it is LightGBM's histogram-binned
design \citep{ke2017lightgbm} implemented inside \code{scikit-learn}
\citep{sklearn2024hgb}, so it stays within the brief's toolset; because it
handles missing values natively, which the cold-start track-record features
need; and because it fits an order of magnitude faster than a comparable
forest, which made a 30-draw search affordable. In e-commerce and logistics the
same family is the working choice -- quantile regression forests for delivery
times at JD.com \citep{salari2022delivery}, tree models for late supplier
deliveries \citep{baryannis2019risk}, and random forests or boosted trees in
the published studies on this very dataset \citep{fouad2023satisfaction,
wangkiat2023satisfaction, gopinadh2025satisfaction}, none of which controls
leakage by prediction moment, the gap Section~\ref{sec:data} closes. The ladder
then supplies what the literature cannot: on this data each rung must beat the
one below it, and Section~\ref{sec:results} reports where it does not.

Every model is a \code{scikit-learn} \code{Pipeline}: median imputation, one-hot
encoding of the four categoricals with rare levels folded into one bucket, then
--- for the linear family only --- standardisation, since trees are invariant to
monotone rescaling and a shared penalty is not \citep{pedregosa2011sklearn}. The
linear family additionally drops one level per categorical. That is a
correction rather than a preference: Phase 1 reported that ordinary least
squares was unusable here, ``rank-deficient, so its coefficients explode'', and
re-running it four ways (Appendix~\ref{app:encoding}) shows the predictions were
always fine (test MAE 4.181, indistinguishable from Ridge) and only the
coefficient vector was not; dropping a level restores full rank, after which
Ridge and OLS agree to three decimals. Regularisation buys numerical stability
here, not accuracy.

\subsection{Cross-validation, and how a model gets chosen}
\label{sec:cv}

Delivery performance drifts hard --- mean lead time falls \textbf{5.34 days}
between training and test --- so shuffled $k$-fold lets a model learn from the
future, the invalidity \citet{bergmeir2012cv} formalise. Hyperparameters are
tuned by \code{TimeSeriesSplit} with five expanding windows inside the training
period, searched with \code{RandomizedSearchCV} \citep{bergstra2012random} over
30 draws per model --- six for the random forests, which cost one to two orders
of magnitude more per fit --- scoring \code{neg\_mean\_absolute\_error} for
regression and \code{average\_precision} for classification.

\emph{Selecting} between tuned models is a separate question, and both obvious
answers fail here for reasons that are properties of the data. Averaging
cross-validated scores fails because the folds differ enormously in difficulty
--- lead time runs 12--17 days across the 2017 folds against 8.3 in the test
window --- so every model's CV error sits about a day above its forward error and
the \emph{ordering} does not survive: cross-validated MAE ranks Ridge above the
boosted trees, while every later window and the test period rank them the other
way round. Using the validation period alone fails too: two months is short, and
these two pull in opposite directions, the late rate being 8.2\% in May 2018 and
\textbf{1.4\%} in June, the lowest in the series.

Each candidate is therefore scored on all six out-of-sample windows ---
the five expanding folds and the validation period --- \textbf{ranked within
each}, and the winner is the best mean rank over the \textbf{three most recent}.
Ranks rather than scores, because ranks are the only thing comparable across
windows of different difficulty and different base rates; three windows because
that covers roughly February to June 2018, recent enough to resemble deployment
and wide enough that neither of the validation period's two odd months decides
alone. All six are reported in \code{report/results/selection\_*.csv}. A
shuffled, customer-disjoint \code{StratifiedGroupKFold} is computed alongside as
a contrast exhibit and never used to choose anything. The selected models are
\emph{not} refitted on train plus validation before the test window is scored: a
production system should, but that would decouple the shipped artifact from the
threshold and calibration fitted against it, and what is reported here is exactly
what is deployed. What the two cross-validation protocols would have reported,
and why their raw PR-AUC is not comparable with the forward window, is
quantified in Appendix~\ref{app:protocol}.

\subsection{Imbalance, metrics aligned to the action, and reproducibility}
\label{sec:imbalance}

Both targets are imbalanced --- 7.4\% late and 11.0\% low review in the test
period --- so accuracy alone is meaningless and every results row carries
\code{majority\_accuracy} beside \code{accuracy}, with PR-AUC reported against
its no-skill floor \citep{saito2015prplot}. Phase 1's literature review argued
from \citet{goorbergh2022imbalance} that correcting imbalance does not improve
discrimination and badly miscalibrates the result, and committed Phase 2 to
testing rather than assuming it (Section~\ref{sec:results-review}). We fit on the
observed prior and move the \emph{decision threshold} instead, tuned on the
validation period by maximising F1 --- which changes the decision without
distorting the probability. Calibrated probabilities are a requirement rather
than a nicety, because Section~\ref{sec:cost} multiplies them by a cost
\citep{elkan2001cost}; they are reported, not assumed
\citep{niculescu2005calibration}. No resampling is used anywhere.

A second alignment follows from Table~\ref{tab:levers}. PR-AUC summarises the
whole ranking and is the right criterion for \emph{choosing} between models,
but neither team acts on the whole ranking: customer experience contacts the
$K$ riskiest deliveries its agents can reach in a day, and operations reviews
the top of the late list. The quantity that describes their day is the precision
and recall of that top-$K$ -- the targeting view of \citet{provost2013data},
and the top-decile lift that profit-driven evaluation uses in place of accuracy
\citep{verbeke2012churn, lessmann2015benchmarking}.
Section~\ref{sec:results-capacity} therefore reports Precision@$K$ and
Recall@$K$ for both flags, with $K$ as a share of the day's volume so that the
same table holds on a quiet Tuesday and on Black Friday, and
Section~\ref{sec:cost} chooses the operating point as a capacity rather than
as a threshold.

\code{SEED = 42} throughout. \code{src/train\_models.py} fits everything and
writes the models, metadata and every result table; \code{src/phase2\_stats.py}
distils those into \code{report/stats\_phase2.txt}, and no number enters this
report that it cannot regenerate. Every figure is drawn from those same tables,
so a figure cannot disagree with one. Artifacts record the \code{scikit-learn}
version that wrote them and the loader checks it;
\code{tests/test\_leakage.py} runs 41 validation checks
(Appendix~\ref{app:repro}).
